% Figure: SI base units and a sample of derived units, showing the
% dependency graph. Each derived unit is shown above the base units it
% combines.
\begin{figure}[htbp]
\centering
\begin{tikzpicture}[
  baseunit/.style={draw=engblack, fill=englime!20, rounded corners=2pt,
    minimum width=2cm, minimum height=0.7cm, font=\small\bfseries},
  derived/.style={draw=engblack, fill=engblue!15, rounded corners=2pt,
    minimum width=2.2cm, minimum height=0.7cm, font=\small},
  arrow/.style={-{Stealth}, thick, draw=engblack!60},
  node distance=0.6cm and 0.4cm,
]
% Base units (bottom row)
\node[baseunit] (m)   at (0, 0)   {metre (m)};
\node[baseunit] (kg)  at (3, 0)   {kilogram (kg)};
\node[baseunit] (s)   at (6, 0)   {second (s)};
\node[baseunit] (A)   at (9, 0)   {ampere (A)};
\node[baseunit] (K)   at (12, 0)  {kelvin (K)};

% Derived units (top row)
\node[derived] (N)    at (1.5, 2.5)  {newton (N)};
\node[derived] (J)    at (4.5, 2.5)  {joule (J)};
\node[derived] (W)    at (7.5, 2.5)  {watt (W)};
\node[derived] (V)    at (10.5, 2.5) {volt (V)};

% N = kg m s^-2
\draw[arrow] (kg) -- (N);
\draw[arrow] (m)  -- (N);
\draw[arrow] (s)  -- (N);
% J = N m = kg m^2 s^-2
\draw[arrow] (N)  -- (J);
\draw[arrow] (m)  -- (J);
% W = J s^-1 = kg m^2 s^-3
\draw[arrow] (J)  -- (W);
\draw[arrow] (s)  -- (W);
% V = W A^-1 = kg m^2 s^-3 A^-1
\draw[arrow] (W)  -- (V);
\draw[arrow] (A)  -- (V);

\end{tikzpicture}
\caption[SI base and derived units]{Five SI base units and four
derived units, with the dependency graph showing how the derived
units are composed. The kilogram, metre, and second carry through
nearly every mechanical unit; the ampere enters at the electrical
boundary. The complete SI maintains seven base units \cite{std:bipm-si2019};
the four further derived units shown here illustrate the
composition pattern.}
\label{fig:vol01:ch01:si-dependency}
\end{figure}
